\section{Directed Graphs and the Cycle Detection Problem}

A \textbf{directed graph} $G = (V, E)$ consists of a finite set of vertices $V$ and a
set of directed edges $E \subseteq V \times V$. In the context of The Round Table Exchange (RTE), each vertex
represents a user and each directed edge $(u, v)$ represents a potential exchange
relationship: user $u$ offers a skill that user $v$ has requested, and their availability
windows overlap within an acceptable geographic range.

A \textbf{directed cycle} of length $k$ is a sequence of distinct vertices
$v_1, v_2, \ldots, v_k$ such that edges $(v_1,v_2), (v_2,v_3), \ldots, (v_{k-1},v_k),
(v_k, v_1)$ all exist in $E$. Finding all such cycles is the core computational problem
this project solves.

\section{Bounded-Depth Depth-First Search}

For small maximum cycle lengths (here $k \leq 4$), a \textbf{bounded-depth
depth-first search (DFS)} is the practical implementation choice. Starting from each
vertex $v$, the algorithm explores outgoing edges recursively to depth $k$, recording
any path that returns to $v$ as a discovered cycle. Visited-node tracking prevents
revisiting within a single path traversal.

Time complexity for exhaustive bounded-DFS over a graph of $n$ nodes and $m$ edges is
$O(n \cdot m^{k-1})$ in the worst case, which is manageable at the target evaluation
scale of up to 1,000 nodes. The algorithm is terminated early per path once depth $k$
is exceeded.

\section{Johnson's Algorithm}

For completeness, \textbf{Johnson's algorithm}~\cite{johnson1975finding} finds all
elementary circuits in a directed graph in time $O((n + m)(c + 1))$, where $c$ is the
number of elementary circuits. It is asymptotically superior to exhaustive DFS for dense
graphs with many cycles, but introduces additional implementation complexity (strongly
connected component decomposition via Tarjan's or Kosaraju's algorithm, a blocked-vertex
stack). This project implements bounded-depth DFS as the primary algorithm given the
constrained maximum cycle length, and notes Johnson's algorithm as a potential upgrade
path if the cycle-length bound is lifted in future work.

\section{Ranking Discovered Cycles}

Once a set of valid cycles is found, cycles are ranked by a composite score:

\[
  \mathrm{score}(C) = \alpha \cdot \mathrm{prox}(C) + (1 - \alpha) \cdot \mathrm{overlap}(C)
\]

where $\mathrm{prox}(C)$ is the normalised mean geographic proximity of participants in
cycle $C$, $\mathrm{overlap}(C)$ is the proportion of shared availability time across all
participants, and $\alpha \in [0,1]$ is a tunable weighting parameter. This two-signal
formula is deliberately minimal: trust weighting and preference scoring are deferred to
future work.

